\setcounter{figure}{0}
\begin{figure}[t!]
\centering
\includegraphics[width=0.86\textwidth]{fig_lanmm_field_p2_twotone.pdf}
\caption{\textbf{LaNMM frequency selection: two-tone TI field on the superficial $P_2$ (gamma)
loop.} The field enters the $P_2$ sigmoid argument as a membrane polarization $\varepsilon$
(mV) with all external drives flat. Color shows the signed change of alpha-band
($8$--$12$~Hz) power from the unstimulated $\varepsilon=0$ baseline, on a zero-centered
scale (arbitrary units). \emph{(a,b)} $P_1$ and $P_2$ enhancement vs.\
$\varepsilon$ and $\Delta f$ (carrier $40$~Hz), peaking at the $P_1$ alpha ($10.0$~Hz,
dashed). The black contour in \emph{(a)} is the frequency-dominance set
$|f_{\rm out}-\Delta f|<0.3$~Hz, drawn because band power
alone cannot separate a driven from an autonomous rhythm; it is a dominance screen, not a
phase-capture test. The $\varepsilon=0$ row is excluded, since with no drive the
autonomous $10.67$~Hz alpha already lies within the tolerance of two grid points (a
resolution floor of $0.5$~Hz); the set widens to $2.75$~Hz at $3$~mV. \emph{(c,d)} Enhancement vs.\ carrier and $\Delta f$ at
$\varepsilon=1.5$~mV: $P_1$ alpha-band enhancement is largest in the
gamma carrier band with $\Delta f$ near the $P_1$ alpha. At $\Delta f=10$~Hz,
the carrier maximum is $44$~Hz. The field-free
$P_2$ gamma peak is near $39.7$~Hz; the finite-drive maximum need not coincide with it.}
\label{fig:lanmm_p2}
\end{figure}
